% Capa
%\begin{titlepage}
% Se quiser uma figura de fundo na capa ative o pacote wallpaper
% e descomente a linha abaixo.
%\ThisCenterWallPaper{0.8}{fsm.jpg}

%\begin{center}
%{\LARGE \nomedoaluno}
%\par
%\vspace{200pt}
%{\Huge \titulo}
%\par
%\vfill
%\textbf{{\large São Paulo}\\
%{\large 2009}}
%\end{center}
%\end{titlepage}

% A partir daqui páginas sem cabeçalho
%\pagestyle{empty}
% Faz com que a página seguinte sempre seja ímpar (insere pg em branco)
%\cleardoublepage

% Números das páginas em algarismos romanos
%\pagenumbering{roman}
\begin{titlepage}
% Página de Rosto

\begin{center}
%\includegraphics[scale=0.1]{fsm.jpg}
%\hspace*{15pt}{\LARGE \nomedolocal}\hspace*{15pt}
\includegraphics[width=\textwidth]{ppgi.jpg}

\par
\vspace{80pt}
{\Huge Um Framework para Medição da Qualidade da Wikipedia com base na Classificação de Surowiecki, J. para Crowdsourcing
}
\par
\vspace{30pt}
{\LARGE Aluno: \Luciana Luciana Sá Brito}
\end{center}
\par
\vspace{90pt}
\hspace*{175pt}\parbox{7.6cm}{{\large Proposta de Dissertação apresentada ao \programa, para a Qualificação ao Título de Mestre em Informática, na Área de Sistemas de Informação.}}

\par
\vspace{1em}
\hspace*{175pt}\parbox{7.6cm}{{\large Orientador: Adriana Vivacqua}}
\par
\vspace{1em}
\hspace*{175pt}\parbox{7.6cm}{\large Co-Orientador: Nome do co-orientador}
\par
\vfill
\begin{center}
\textbf{{\large \localidade}\\
{\large \the\year}}
\end{center}
\end{titlepage}
%\newpage

% Ficha Catalográfica
%\hspace{8em}\fbox{\begin{minipage}{10cm}
%Aluno, Nome C.

%\hspace{2em}\titulo

%\hspace{2em}\pageref{LastPage} páginas

%\hspace{2em}Dissertação (Mestrado) - Instituto de Biociências da Universidade de São Paulo. Departamento de XXXXXXXX.

%\begin{enumerate}
%\item Palavra-chave
%\item Palavra-chave
%\item Palavra-chave
%\end{enumerate}
%I. Universidade de São Paulo. Instituto de Biociências. Departamento de XXXXXXXX.

%\end{minipage}}
%\par
%\vspace{2em}
%\begin{center}
%{\LARGE\textbf{Comissão Julgadora:}}

%\par
%\vspace{10em}
%\begin{tabular*}{\textwidth}{@{\extracolsep{\fill}}l l}
%\rule{16em}{1px} 	& \rule{16em}{1px} \\
%Prof. Dr. 		& Prof. Dr. \\
%Nome			& Nome
%\end{tabular*}

%\par
%\vspace{10em}

%\parbox{16em}{\rule{16em}{1px} \\
%Prof. Dr. \\
%Nome do Orientador}
%\end{center}

%\newpage

% Dedicatória
% Posição do texto na página
%\vspace*{0.75\textheight}
%\begin{flushright}
%  \emph{Dedicatória...}
%\end{flushright}

%\newpage

% Epígrafe
%\vspace*{0.4\textheight}
%\noindent{\LARGE\textbf{Exemplo de epígrafe}}
% Tudo que você escreve no verbatim é renderizado literalmente (comandos não são interpretados e os espaços são respeitados)
%\begin{verbatim}
%O que é bonito?
%É o que persegue o infinito;

%\end{verbatim}
%\begin{flushright}
%Lenine e Bráulio Tavares
%\end{flushright}

%\newpage

% Agradecimentos

% Espaçamento duplo
%\doublespacing

%\noindent{\LARGE\textbf{Agradecimentos}}

%Agradeço ao meu orientador, ao meu co-orientador, aos meus colaboradores, aos técnicos, à seção administrativa, à fundação que liberou verba para minhas pesquisas, aos meus amigos, à minha família e ao meu grande amor.

\newpage
\pagenumbering{roman}
\vspace*{10pt}
% Abstract
\begin{center}
  \emph{\begin{large}Resumo\end{large}}\label{resumo}
\vspace{2pt}
\end{center}
% Pode parecer estranho, mas colocar uma frase por linha ajuda a organizar e reescrever o texto quando necessário.
% Além disso, ajuda se você estiver comparando versões diferentes do mesmo texto.
% Para separar parágrafos utilize uma linha em branco.
\noindent
colocar resumo aqui
\par
\vspace{1em}
\noindent\textbf{Palavras-chave:} \tagstrabalho
%\newpage

% Criei a página do abstract na mão, por isso tem bem mais comandos do que o resumo acima, apesar de serem idênticas.
%\vspace*{10pt}
% Abstract
%\begin{center}
%  \emph{\begin{large}Abstract\end{large}}\label{abstract}
%\vspace{2pt}
%\end{center}

% Selecionar a linguagem acerta os padrões de hifenação diferentes entre inglês e português.
%\selectlanguage{english}
%\noindent
%This is the most important part of your work.
%This is what most people will read.
%Be concise without omitting content.
%A good abstract explains why this is an interesting study, tells how it was done, what was found, contextualizes the results and set conclusions.
%\par
%\vspace{1em}
%\noindent\textbf{Keywords:} word1, word2, word3

% Voltando ao português...
%\selectlanguage{brazilian}

%\newpage

% Lista de figuras
%\listoffigures

% Lista de tabelas
%\listoftables

% Abreviações
% Para imprimir as abreviações siga as instruções em 
% http://code.google.com/p/mestre-em-latex/wiki/ListaDeAbreviaturas
%\printnomenclature

% Índice

\tableofcontents

